% !TEX encoding = UTF-8
% !TEX program = lualatex

\chapter{Теоретические основы атрибутивных метаграфов и их применение в кибербезопасности}

\section{Формальное определение метаграфа}

Метаграф представляет собой обобщение гиперграфа, позволяющее моделировать сложные зависимости, в которых как источники, так и цели действия могут быть множествами объектов. Впервые строгое определение метаграфа было предложено в работах Zhang и Thulasiraman~\cite{Zhang2006}, однако его применение в области кибербезопасности остаётся малоизученным.

Согласно~\cite{Latypov2020}, метаграф эффективно описывает многоуровневые модели компьютерных атак, где каждый уровень соответствует определённой тактике или технике. В настоящей работе используется следующее формальное определение.

\begin{definition}
Метаграфом называется упорядоченная четвёрка
\[
G = (V, E, \mathrm{src}, \mathrm{tgt}),
\]
где:
\begin{itemize} 
    \item $V$ — конечное непустое множество вершин (объектов системы);
    \item $E$ — конечное непустое множество гипердуг (событий или действий);
    \item $\mathrm{src}: E \rightarrow 2^V$ — функция, сопоставляющая каждой дуге множество источников;
    \item $\mathrm{tgt}: E \rightarrow 2^V$ — функция, сопоставляющая каждой дуге множество целей.
\end{itemize}   
    При этом для любой дуги $e \in E$ выполняется условие $\mathrm{src}(e) \cap \mathrm{tgt}(e) = \varnothing$.
\end{definition}

В контексте кибербезопасности вершины могут представлять процессы, файлы, пользователей, сетевые соединения, а гипердуги — события безопасности, такие как «процесс~A создал файл~B и установил соединение с хостом~C».

\section{Атрибутивные метаграфы}

Для адекватного моделирования компьютерных атак необходимо расширить метаграф за счёт семантических и количественных характеристик. Следуя подходу, предложенному в~\cite{Latypov2020}, вводится понятие атрибутивного метаграфа.

\begin{definition}
Атрибутивным метаграфом называется шестёрка
\[
G_{\mathrm{attr}} = (V, E, \mathrm{src}, \mathrm{tgt}, A_V, A_E, \alpha, \beta),
\]
где:
    \item $A_V$ — множество атрибутов вершин;
    \item $A_E$ — множество атрибутов дуг;
    \item $\alpha: V \rightarrow A_V$ — функция атрибутизации вершин;
    \item $\beta: E \rightarrow A_E$ — функция атрибутизации дуг.
\end{definition}

Типичные атрибуты вершин включают:
\begin{itemize}
    \item \texttt{type} $\in \{\texttt{process}, \texttt{file}, \texttt{user}, \texttt{ip}, \texttt{registry\_key}\}$;
    \item \texttt{name} — идентификатор объекта;
    \item \texttt{hash} — криптографический хеш (для файлов);
    \item \texttt{integrity\_level} — уровень целостности (Low, Medium, High, System).
\end{itemize}
Атрибуты дуг включают:
\begin{itemize}
    \item \texttt{mitre\_ttp} — идентификатор техники MITRE ATT\&CK (например, \texttt{T1059.001});
    \item \texttt{timestamp} — временная метка события;
    \item \texttt{confidence} — достоверность события (0.0–1.0);
    \item \texttt{source} — источник данных (EDR, Zeek, Windows Event Log).
\end{itemize}

Такой подход позволяет интегрировать модель с фреймворком MITRE ATT\&CK, что обеспечивает интерпретируемость и совместимость с открытыми источниками угроз, включая описания APT-группировок из США, Великобритании и Израиля (например, Lazarus, APT28, Charming Kitten), при условии использования публичных отчётов~\cite{Mandiant2023,CrowdStrike2024}.

\section{Семантическое расширение для моделирования атак}

Каждая дуга $e$ с атрибутом $\beta(e).\texttt{mitre\_ttp} = t$ автоматически наследует тактику $T(t)$ из фреймворка MITRE ATT\&CK. Это позволяет строить тактические профили метаграфов.

\begin{definition}
Тактическим профилем метаграфа называется отображение
\[
\tau: \mathcal{T} \rightarrow 2^E,
\]
где $\mathcal{T}$ — множество тактик MITRE ATT\&CK, и
\[
\tau(T) = \{ e \in E \mid T(\beta(e).\texttt{mitre\_ttp}) = T \}.
\]
\end{definition}

Например, для APT29 (Cozy Bear) тактический профиль включает тактики \textit{Initial Access} (T1566), \textit{Execution} (T1059), \textit{Persistence} (T1547) и \textit{Lateral Movement} (T1021).

\section{Алгебра метаграфов}

Для манипуляции метаграфами вводятся базовые операции, адаптированные из~\cite{Zhang2006} и расширенные в~\cite{Latypov2020}.

\subsection{Объединение}

Пусть $G_1$ и $G_2$ — два атрибутивных метаграфа. Их объединение определяется как
\[
G = G_1 \cup G_2 = (V_1 \cup V_2, E_1 \cup E_2, \mathrm{src}, \mathrm{tgt}, \alpha, \beta),
\]
где функции $\mathrm{src}$, $\mathrm{tgt}$, $\alpha$, $\beta$ определены естественным образом на объединении.

\subsection{Композиция}

Композиция $G = G_1 \circ G_2$ моделирует последовательное выполнение двух этапов атаки и определена, если множества целей $G_1$ и источников $G_2$ пересекаются.

\subsection{Проекция}

Проекция по тактике $T$:
\[
\pi_T(G) = (V', E', \mathrm{src}', \mathrm{tgt}', \alpha', \beta'),
\]
где $E' = \tau(T)$, $V' = \bigcup_{e \in E'} (\mathrm{src}(e) \cup \mathrm{tgt}(e))$.

Аналогично определяется проекция по временному интервалу $[t_1, t_2]$.

\section{Формализация атак с помощью метаграфов}

Рассмотрим формализацию типичной фишинговой атаки, характерной для APT-группировок~\cite{Latypov2020}. Этапы:
\begin{enumerate}
    \item Доставка: файл \texttt{invoice.docm} (T1566);
    \item Выполнение: запуск макроса → \texttt{powershell.exe} (T1059.001);
    \item Устойчивость: создание задачи в Планировщике (T1053.005);
    \item C2: HTTPS-соединение с сервером (T1071.001).
\end{enumerate}

Соответствующий метаграф $G_{\mathrm{phish}}$ строится по определению~2.2. Такая формализация позволяет использовать метаграф как шаблон для последующего сопоставления.

\section{Сравнение с альтернативными моделями}

В работах~\cite{Novokhrestov2014,Zegzhda2016,Kotenko2013,Gorbachev2018} рассматриваются различные подходы к моделированию угроз: от многоагентных систем до многоуровневых архитектур защиты. Однако ни одна из моделей не обеспечивает одновременно:
\begin{itemize} 
\item поддержку множественных источников и целей;
    \item интеграцию с MITRE ATT\&CK;
    \item динамическое обновление по потоковым данным;
    \item вычислительную эффективность сопоставления.
\end{itemize}
Атрибутивный метаграф, как показано в~\cite{Latypov2020}, обеспечивает оптимальный баланс между выразительной мощностью и практической применимостью, что подтверждается экспериментами в главе~4.

\section{Выводы по главе}

В настоящей главе:
\begin{itemize}
    \item определены операции алгебры метаграфов;
    \item показана интеграция с MITRE ATT\&CK и возможность моделирования APT-атак;
    \item обосновано преимущество предложенной модели перед альтернативами.
\end{itemize}   
Полученные теоретические результаты лежат в основе методики выявления атак, описанной в главе~3.